\section{Temperature、Sampling 与 Oracle Approximation}

Sampling strategy 决定从 model distribution 的哪个区域取候选。\term{temperature}（温度）对 logits $z_i$ 做
\[
p_T(i)=\frac{\exp(z_i/T)}{\sum_j\exp(z_j/T)}.
\]
$T<1$ 使分布更尖、偏向高概率 token；$T>1$ 提高多样性，也增加语法/语义错误。最佳温度取决于 metric 和 $k$。

\subsection{为 pass@k 校准 temperature}

本节先把 decoding 参数与评价预算连接起来。温度不应凭习惯固定，而要根据 $k$、verifier 和延迟预算校准。低温让候选集中在 mode，适合一次展示；高温扩大 search coverage，但只有可靠 tests 能把 diversity 转成最终正确性。

\lecturefigure{slide-39-temperature-calibration.jpg}{低温有利 pass@1，较高温可能为 large-k 提供多样候选。}{官方视频 41:26。}

\begin{knowledgebox}{读图：每条曲线对应一个搜索预算}
当 $k=1$，过高温度会把质量分散到低概率错误；当 $k$ 较大，适度 diversity 可避免重复同一错误模式。最优温度不是模型常数，应在 validation tasks 上按实际 generation budget 校准，并防止 benchmark overfitting。
\end{knowledgebox}

\subsection{Sampling 的“不合理有效性”}

接下来观察增加 candidate budget 的收益曲线，并把 benchmark gain 与推理、执行、筛选成本放在同一决策中。曲线的横向移动意味着更多 samples，纵向提升则是更多题至少出现一个正确候选；二者之间不是免费的能力增长。

\lecturefigure{slide-40-sampling-effectiveness.jpg}{随着 samples 增加，pass@k 可远高于单次输出表现。}{官方视频 42:19。}

\begin{knowledgebox}{读图：提升来自 search，不一定来自更好单样本}
曲线说明同一模型在更大 candidate budget 下可解决更多题。真实系统还要支付 inference、sandbox、test 与 ranking 成本。若 tests 不完备，选择器可能偏向 exploit；如果多个候选高度相关，有效样本数低于名义 $k$。
\end{knowledgebox}

\teachervoice{讲者把 sampling 称为“unreasonably effective”。课堂提示不是鼓励盲目多采样，而是提醒模型质量是一整个 conditional distribution；产品只展示一个 completion 时，decoding 和 verifier 共同决定可见能力。}

\subsection{近似 oracle}

进一步引入选择上限。\term{oracle} 在这里是假想选择器：只要 $k$ 个候选中存在正确程序，就能识别并返回它。Unit tests 是有限 oracle approximation，覆盖越完整，选择越可靠。这个抽象把“模型没生成正确答案”和“系统没认出正确答案”分成两个可测瓶颈。

\lecturefigure{slide-41-oracle-sampling.jpg}{Sampling 加 verifier 近似“从候选集合中选出正确程序”的 oracle。}{官方视频 43:17。}

\begin{knowledgebox}{读图：比较 model curve 与 oracle curve 的差距}
若 oracle-style pass@k 很高而实际 selector 低，瓶颈在 verification/ranking；若两者都低，瓶颈在 model distribution。工程团队应分别投资 generation、tests、static analysis 与 reranking，不能把所有差距归因于模型参数。
\end{knowledgebox}

\begin{warningbox}{测试可以被满足，也可以被投机}
有限 tests 允许 overfit、hard-code 或触发未覆盖副作用。生成代码必须在 sandbox 运行，并检查 timeout、memory、network、filesystem 与 dependency。Oracle 是抽象上限，不是现实安全保证。
\end{warningbox}

\subsection{Codex-S：监督信号继续塑形}

Codex-S 在 code-pretrained model 上使用 stand-alone Python functions 与正确 solutions 做 supervised fine-tuning。它改变了 prompt-to-solution distribution，使 HumanEval-style tasks 更匹配。

\lecturefigure{slide-42-codex-s-training.jpg}{Codex-S 用监督函数数据进一步对齐 HumanEval-style completion。}{官方视频 44:10；Chen et al. 2021。}

\begin{knowledgebox}{读图：Fine-tuning 改善 task fit，也可能缩窄分布}
Supervised examples 明确教模型从 docstring 到 function body，提高 benchmark performance。需要检查是否牺牲其他语言、repository context 或 diversity，并防止 train-test similarity。Fine-tuning gain 不代表 base model 不重要，而是 objective 更接近 evaluation。
\end{knowledgebox}

\subsection{Main figure：模型规模、Codex 与 sampling budget}

本节综合 model family、scale 与 sampling budget，区分 pretraining gain、code fine-tuning gain 和 search gain，避免把三种效果混成一个数字。

\lecturefigure{slide-43-main-figure.jpg}{Codex main figure 联合比较 model family、规模与 pass@k。}{官方视频 45:52；Chen et al. 2021。}

\begin{knowledgebox}{读图：先读图例，再读斜率和间距}
横轴/曲线通常编码 model size 或 sampling budget；不同颜色区分 GPT、Codex、Codex-S。比较同一 $k$ 下 fine-tuning gain，再比较同一模型从 pass@1 到 pass@100 的 search gain。不要把 pass@100 当交互式单 completion 体验，也不要忽略生成 100 个候选的成本。
\end{knowledgebox}

\subsection{本章小结}

Temperature 控制质量—多样性，pass@k 把候选预算纳入评价，tests 近似 oracle，Codex-S 则通过 supervised data 改变分布。代码系统的能力来自 model、decoding 与 verifier 三者。

